% Step 0: the naive baseline. 14 frames, DEMO not lab, session 0:20-0:35.
\section{Step 0. The naive baseline}\subsection{}

% Step 0 opener. Step table, this row tinted.
\begin{frame}{Step 0. The naive baseline}
\begin{center}
\tabfont
\begin{tabular}{cll}
\toprule
Step & What we build & The concept it forces \\
\midrule
\rowcolor{yo!10}\yob{0} & \yob{Naive baseline: one prompt} & \yob{A passing check is only as good as its scope} \\
1 & Write a spec, then build & \ybb{Plan mode}, \ybb{sub-agents} for independent verification \\
2 & Generalize across quarters & \ybb{Skills}, and the tie-out as a test suite \\
3 & Analysis and report & Analysis as verification; \ybb{grounding} in your references \\
\bottomrule
\end{tabular}
\end{center}
 
\vs
\begin{itemize}
    \item I will walk through my own output in detail; you do not need to type anything.
    \item It took $\sim$15 to 20 minutes to finish this one simple prompt.
    \item I experimented with three different models: Opus 5 at xhigh effort, Opus 5 at medium effort, and Sonnet 5 at medium effort.
\end{itemize}
\end{frame}

% The prompt, verbatim.
\begin{frame}[fragile]{Start with the laziest possible prompt}
\begin{prompt}
Find a large publicly traded BDC, download its most recent 10-Q from SEC EDGAR, and parse the financial statements and the schedule of investments into two structural tables I can analyze.
\end{prompt}
\vs
\begin{wideitemize}
  \item Not in it: a ticker, a schema, a file format, a variable list.
  \item Not in it: how to download or paste the filing.
  \item Not in it: any word about verifying anything.
\end{wideitemize}
\end{frame}

% It picked its own target and surveyed the filing before parsing it.
% MERGED 2026-08-19 (clock): was two frames, 'It picked its own target' and
% 'It reads the document before parsing it'.  Split them again if the block
% ever gets its time back; both screenshots are still here.
\begin{frame}{It picked its own target, and looked before it parsed}
\begin{columns}
\begin{column}{0.5\textwidth}
\begin{center}
\includegraphics[width=0.9\textwidth]{bdc_design_pharse0-1.png}
\end{center}
\end{column}
\begin{column}{0.5\textwidth}
\begin{center}
\includegraphics[width=0.9\textwidth]{bdc_design_pharse0-2.png}
\end{center}
\end{column}
\end{columns}
\vs
\begin{itemize}\small
  \item \ybb{Ares Capital (ARCC)}, reasoned from ``large publicly traded BDC'',
        and its \ybb{most recent 10-Q} (2026-06-30): a 24.8 MB inline-XBRL document.
  \item It then \ybb{surveyed the filing} before parsing -- how many tables, where
        the schedule starts, how the columns repeat.
  \item ``\ybb{Look before you act}'' costs a couple of minutes, and it is why the parse works.
\end{itemize}
\end{frame}

% What one sentence produced.
\begin{frame}{What came back, from one sentence}
\begin{columns}
\begin{column}{0.5\textwidth}
\begin{center}
\includegraphics[width=\textwidth]{bdc_design_pharse0-10.png}
\end{center}
\end{column}
\begin{column}{0.5\textwidth}
\begin{itemize}
    \item It took $\sim$22 minutes. 
    \item Two tables: 435 statement facts; \ybb{2,847} investment rows
    (1,439 at 2026-06-30, 1,408 at 2025-12-31).
    \item Plus a footnote legend, a subtotals file, a README, and a
    \ybb{re-runnable Python parser} -- 803 lines of it.
\end{itemize}
\end{column}
\end{columns}
\end{frame}

% The 948 percent coupon: the filing is wrong, not the agent.
\begin{frame}[fragile]{948.00\% coupon -- and the filing is wrong, not the agent}
\begin{columns}
\begin{column}{0.4\textwidth}
\begin{center}
\includegraphics[width=\textwidth]{bdc_design_pharse0-7.png}
\end{center}
\end{column}
\begin{column}{0.6\textwidth}
\begin{itemize}
  \item \yob{It also found an obvious bug}: Adonis Bidco shows a coupon of 948.00\%.
  \item Then it dug in: three of the 1,846 coupon facts (Adonis Bidco, Covert HoldCo, Mountaineer Merger) are tagged \file{scale="0"} where every other rate carries \file{scale="-2"}. Taking the XBRL value at face value turns a 9.48\% coupon into 948\%.
  \item The mis-tagging is \yob{ARCC's, sitting live at SEC}. 
  \item The agent noticed it, worked around it, and wrote it down.
\end{itemize}
\end{column}
\end{columns}
\end{frame}

% It invented a tie-out nobody asked for.
\begin{frame}{Nobody asked for verification. It built one.}
\begin{center}
\tabfont
\begin{tabular}{p{8cm}ll}
\toprule
Check, each against a number the filing prints & 2026-06-30 & 2025-12-31 \\
\midrule
A. parsed lines vs printed Total Investments & exact & exact \\
B. each issuer subtotal vs the lines above it & 394/394 agree & 380/380 agree \\
C. schedule total vs balance sheet & within \$0.4m & within \$0.2m \\
D. affiliation split vs the three balance sheet buckets & within \$0.8m & within \$0.8m \\
E. investments / equity vs the printed \% of net assets & 211.28 vs 211.32 & 205.93 vs 205.93 \\
\bottomrule
\end{tabular}
\end{center}
\vs
The word ``verify'' appears \ybb{nowhere} in the prompt, and it invented
a balance-sheet tie-out.
\vs
If we stopped here, we would think the lesson was ``the prompt can be lazy''.
\end{frame}

% The three-run experiment: same prompt, three configurations.
\begin{frame}{Same prompt, three times}
\begin{wideitemize}
  \item Identical prompt, identical filing, three models:
  \begin{itemize}
    \item Opus 5 at xhigh effort 
    \item Opus 5 at medium effort (max > xhigh > high > medium > low)
    \item Sonnet 5 at medium effort (Fable 5 > Opus 5 > Sonnet 5)
  \end{itemize}
  \item Nothing else differs. No extra instruction, no follow-up, no correction.
  \item This is the experiment sixty of you run every time you all point the same
        tool at the same task: \yob{same instruction, different machine and model.}
\end{wideitemize}
\takeaway{Question: what do you expect to differ -- and what do you expect to stay the same?}
\end{frame}

% Three schemas that cannot be pooled.
\begin{frame}{Three right answers with no shared vocabulary}
All three chose ARCC, parsed both periods, and passed a self-invented tie-out at 0.001\%. \ybb{Every run succeeded.} Then:
\vs
\tabfont
\begin{center}
\begin{tabular}{llll}
\toprule
 & Opus 5 xhigh & Opus 5 medium & Sonnet 5 medium \\
\midrule
Columns & 44 & 22 & 26 \\
Company field & \file{portfolio\_company} & \file{issuer} & \file{issuer} \\
Fair value field & \file{fair\_value\_usd} & \file{fair\_value} & \file{fair\_value\_usd} \\
Type field & \file{investment\_type} & \file{instrument} & \file{investment\_type} \\
Verification shipped as & \file{tieout\_report.txt} & separate script & inside the parser \\
Parser & 803 lines, 1 file & 526 lines, 2 files & 504 lines, 2 files \\
\bottomrule
\end{tabular}
\end{center}
\takeaway{There are sixty of you in this room. That is sixty correct, but incompatible datasets.}
\end{frame}

% Row counts, and what the 1,409 means.
\begin{frame}{Count the rows. What does the 1,409 mean?}
\tabfont
\begin{center}
\begin{tabular}{llrrr}
\toprule
Run & date & rows & distinct companies & investment type is empty? \\
\midrule
Opus 5 xhigh & 2026-06-30 & 1,439 & 593 & 0 \\
Opus 5 xhigh & 2025-12-31 & 1,408 & 580 & 0 \\
Opus 5 medium & 2026-06-30 & 1,439 & 605 & 8 \\
\yob{Opus 5 medium} & \yob{2025-12-31} & \yob{1,409} & \yob{1,409} & \yob{1,409} \\
Sonnet 5 medium & 2026-06-30 & 1,439 & 593 & 0 \\
Sonnet 5 medium & 2025-12-31 & 1,409 & 580 & 1 \\
\bottomrule
\end{tabular}
\end{center}
\vs
\takeaway{Why does Opus 5 medium have 1,409 distinct companies, but all of them have an empty investment type field?}
\end{frame}

% The failure, in the data.
\begin{frame}[fragile]{Every row became its own borrower}
\ybb{Opus 5 medium}, at 2025-12-31:
\begin{shellout}
issuer                                                                               investment_type
Automotive Keys Group, LLC and Automotive Keys Investor, LLC, Class A common units              -
Automotive Keys Group, LLC and Automotive Keys Investor, LLC, First lien senior secured loan    -
\end{shellout}
\ybb{Opus 5 xhigh}, the same rows:
\begin{shellout}
portfolio_company                                                  investment_type
Automotive Keys Group, LLC and Automotive Keys Investor, LLC       Class A common units
Automotive Keys Group, LLC and Automotive Keys Investor, LLC       First lien senior secured loan
\end{shellout}
\begin{itemize}
    \item The problem: the instrument type text has been merged into the borrower name. It parsed the current period correctly and the comparative period wrongly.
    \item \yob{And it ties out to the balance sheet exactly, on both dates.} The check was measuring \ybb{dollars} while the thing that broke was \yob{structure}. 
\end{itemize}
\end{frame}

% Effort, not model tier. Closes Step 0.
\begin{frame}{The broken one is Opus. Sonnet at the same effort is clean.}
\begin{wideitemize}
  \item \ybb{Sonnet 5 at medium effort is clean}: 593 and 580 distinct companies.
  \item Model tier and reasoning effort \yob{did not predict correctness} here.
  \item What the extra effort bought was \ybb{coverage}, not correctness: 44 columns
        vs 26, 14 decoded footnote flags, a subtotals file, a legend file.
\end{wideitemize}
\takeaway{``Use the biggest model'' is not the lesson. How to design verification is the one.}
\end{frame}
